% Part: normal-modal-logic
% Chapter: completeness
% Section: canonical-models

\documentclass[../../../include/open-logic-section]{subfiles}

\begin{document}

\olfileid{nml}{com}{cmd}

\olsection{కానానికల్ నమూనాలు}

మోడల్ వ్యవస్థ~$\Sigma$కు \emph{కానానికల్ నమూనా} అంటే
ఒక నిర్దిష్ట నమూనా~$\mModel{M^\Sigma}$. సంపూర్ణ
$\Sigma$-అవిరుద్ధ సమితులన్నీ దానిలో లోకాలు. ఒక లోకం~$\Delta$లో
సత్యమయ్యే \tetoken{సూత్రాలు}{!!{formula}s} సరిగ్గా~$\Delta$లో
ఉన్న \tetoken{సూత్రాలే}{!!{formula}s} అయ్యేలా
ప్రాప్యత సంబంధం~$R^\Sigma$నూ, మూల్యనిర్ణయం $V^\Sigma$నూ నిర్వచిస్తాం.

\begin{defn}
  $\Sigma$ నార్మల్ మోడల్ తర్కం అనుకుందాం. $\Sigma$కు చెందిన
  \emph{కానానికల్ నమూనా} $\mModel{M}^\Sigma = \tuple{W^\Sigma,
    R^\Sigma, V^\Sigma}$.
  కానానికల్ నిర్మాణాన్ని నమూనాగా వాడేటప్పుడు వ్యవస్థ అవిరుద్ధమని తీసుకుంటాం.
  అందువల్ల లోకాల సమితి ఖాళీ కాదు.
  ఇక్కడ:
  \begin{enumerate}
  \item $W^\Sigma = \Setabs{ \Delta }{\Delta \text{
      సంపూర్ణ $\Sigma$-అవిరుద్ధం} }$.
  \item $R^\Sigma \Delta\Delta'$ అయ్యేది
    \iftag{prvBox}{$\Box^{-1}\Delta \subseteq
      \Delta'$}{$\Diamond\Delta' \subseteq \Delta$} అయినప్పుడు, అప్పుడే మాత్రమే.
  \item $V^\Sigma(p) = \Setabs{\Delta}{p \in \Delta}$.
  \end{enumerate}
\end{defn}


\end{document}
